\chapter{Th\'eorie de la sp\'ecialisation du~groupe~fondamental}
\label{X}
Dans le pr\'esent expos\'e, nous nous bornons \`a l'\'etude du
groupe fondamental des fibres g\'eom\'etriques dans un morphisme
\emph{propre}, \ie du groupe fondamental d'un sch\'ema
alg\'ebrique \emph{propre} variable. Dans un expos\'e
ult\'erieur, nous g\'en\'eraliserons la technique employ\'ee
aux rev\^etements \'etales \emph{modérément ramifiés}
\og \`a l'infini\fg. Cela nous donnera par exemple une solution du
\og Probl\`eme des trois points\fg dans le cas des rev\^etements
galoisiens d'ordre premier \`a la caract\'eristique (\ie une
d\'etermination des rev\^etements galoisiens de la droite
$\PP_k^1$, ramifi\'es au plus en trois points donn\'es et
mod\'er\'ement ramifi\'es en ces points), et de ses variantes
\'evidentes.

\section[La suite exacte d'homotopie pour un morphisme propre]
{La suite exacte d'homotopie pour un morphisme propre et s\'eparable}
\label{X.1}

\begin{definition}
\label{X.1.1}
Un pr\'esch\'ema $X$ sur un corps $k$ est dit
\emph{séparable}, ou \emph{séparable sur~$k$}, si pour toute
extension $K$ de~$k$, $X\otimes_k K$ est r\'eduit. Si $f\colon X\to
Y$ est un morphisme de pr\'esch\'emas, on dit que $f$ est
\emph{séparable}, ou que $X$ \emph{est séparable sur~$Y$}, si
$X$ est plat sur~$Y$ et si pour tout $y\in Y$, la fibre $X\otimes_Y
\kres(y)$ est s\'eparable sur~$\kres(y)$.
\end{definition}

Si $X$ est un pr\'esch\'ema sur un corps~$k$, dire qu'il est
s\'eparable signifie aussi qu'il est \emph{réduit}, et que les
corps~$\kres(x)$, pour $x$ point g\'en\'erique d'une composante
irr\'eductible de~$X$, sont des extensions s\'eparables de~$k$. Si
$k$ est parfait, il revient donc au m\^eme de dire que $X$ est
r\'eduit. Notons que si $X$ est s\'eparable sur~$Y$, alors pour
tout changement de base $Y'\to Y$, $X'=X\times_Y Y'$ est s\'eparable
sur~$Y'$. On peut prouver aussi, moyennant des hypoth\`eses de
finitude convenables, que le compos\'e de deux morphismes
s\'eparables est
un morphisme s\'eparable. Nous en aurons besoin seulement sous la
forme suivante: \emph{si $X$ est séparable sur~$Y$, et $X'$
étale sur~$X$, $X'$ est séparable sur~$Y$.} C'est en effet
une cons\'equence imm\'ediate des d\'efinitions
et~I~\ref{I.9.2}. Par ailleurs, l'hypoth\`ese \og morphisme
s\'eparable\fg nous servira par l'interm\'ediaire de la proposition
suivante:
\begin{proposition}
\label{X.1.2}
Soit $f\colon X\to Y$ un morphisme propre et s\'eparable, avec $Y$
localement noeth\'erien, et consid\'erons sa factorisation de
\emph{Stein} $X\lto{f'}Y'\dpl\to Y$, (o\`u
$f_*'(\cal{O}_X)=\cal{O}_{Y'}$, $Y'$ \'etant fini sur~$Y$ et
isomorphe au spectre de l'alg\`ebre $f_*(\cal{O}_X)$). Alors $Y'$
est un \emph{revêtement étale} de~$Y$.
\end{proposition}

Cette proposition figurera dans (EGA~III 7)\footnote{\Cf EGA~III
7.8.10 (i)}. Indiquons le principe de la d\'emonstration. On se
ram\`ene facilement au cas o\`u $Y$ est le spectre d'un anneau
local complet~$A$, et faisant encore une extension finie plate
convenable de ce dernier (correspondant \`a une extension
r\'esiduelle convenable), on peut supposer que les composantes
connexes de la fibre du point ferm\'e $y$ sont
g\'eom\'etriquement connexes, ce qui signifie aussi que
$\H^0(X_y,\cal{O}_{X_y})$ se d\'ecompose en un produit de corps
identiques \`a $k=\kres(y)$. Supposant alors $X$ connexe, ce qui est
loisible, on aura $\H^0(X_y,\cal{O}_{X_y})=k$, donc l'homomorphisme
$A\to \H^0(X_y,\cal{O}_{X_y})$ est \emph{surjectif}. On en conclut par
une proposition g\'en\'erale (du type K\"unneth) que
$f_*(\cal{O}_X)$ est d\'efini par un module $B$ sur~$A$ qui est
libre sur~$A$, et que $B/\goth{m} B\to \H^0(X_y,\cal{O}_{X_y})=k$ est
bijectif. Donc en l'occurrence $B$ est une alg\`ebre \'etale
sur~$A$, ce qui ach\`eve la d\'emonstration.

\begin{theoreme}
\label{X.1.3}
Soit $f\colon X\to Y$ un morphisme propre et s\'eparable, avec $Y$
localement noeth\'erien et connexe, et supposons
$f_*(\cal{O}_X)=\cal{O}_Y$ (ce qui implique que les fibres de~$X$ sur~$Y$ sont g\'eom\'etriquement connexes, et r\'eciproquement
gr\^ace \`a~\ref{X.1.2}). Soient~$y$ un point de~$Y$,
$\overline{\kres(y)}$ une cl\^oture alg\'ebrique de~$\kres(y)$,
$\overline X_y =X_y\otimes_{\kres(y)} \overline{\kres(y)}$. Soient \hbox{enfin~$X'$} un rev\^etement \'etale \emph{connexe} de~$X$, et $\overline
X_y' =X_y'\otimes_{\kres(y)} \overline{\kres(y)}$. Pour qu'il existe un
rev\^etement \'etale $Y'$ de~$Y$ et un $X$-isomorphisme
$X'\isomto X\times_Y Y'$, il faut et il suffit que $\overline{X'}_y$
admette une section sur~$\overline X_y$.
\end{theoreme}

Posant $Y'=\Spec(h_*(\cal{O}_{X'}))$ (o\`u $h\colon X'\to Y$ est le
morphisme compos\'e $X'\to X\to Y$), il suffit de prouver que le
$Y$-morphisme canonique
$$
X'\to X\times_Y Y'
$$
est un \emph{isomorphisme}, et que $Y'$ est \'etale sur~$Y$. Or nous
savons d\'ej\`a par~\ref{X.1.2} que $Y'$ est \'etale sur~$Y$,
donc $X\times_Y Y'$ est \'etale sur~$X$, donc le morphisme $X'\to
X\times_Y Y'$ est \'egalement \'etale (I~\ref{I.4.8}). D'ailleurs,
$Y'$ est connexe comme image de~$X'$ qui l'est, donc $X\times_Y Y'$
est connexe puisque $X$ est \`a fibres connexes sur~$Y$
(IX~\ref{IX.3.4} et V~\ref{V.6.9}~(iii)). Donc pour prouver que
$X'\to X\times_Y Y'$ est un isomorphisme, il suffit de voir que son
degr\'e de projection en \emph{un} point de~$X\times_Y Y'$ est
\'egal \`a~1. Or ceci r\'esulte facilement de l'hypoth\`ese
que $\overline{X'}_y$ admet une section sur~$\overline X_y$, soit
par utilisation de IX~\ref{IX.6.6}, soit plus simplement en notant qu'il
suffit de prouver l'existence d'un tel point dans $X\times_Y Y'$
apr\`es changement de base
$\Spec(\overline{\kres(y)})\to Y$, o\`u cela
est \'evident. Cela ach\`eve la d\'emonstration de~\ref{X.1.3}.

Tenant compte de IX~\ref{IX.3.4} et du dictionnaire V~\ref{V.6.9} et
V~\ref{V.6.11}, on peut mettre~\ref{X.1.3} sous la forme
\'equivalente suivante

\begin{corollaire}
\label{X.1.4}
Avec les notations pr\'ec\'edentes pour $f\colon X\to Y$ et
$\overline X_y$, soit $\bar{a}$ un point g\'eom\'etrique de
$\overline X_y$, $a$ son image dans $X$ et $b$ son image
dans~$Y$. Alors la suite suivante d'homomorphismes de groupes est
\emph{exacte}:
$$
\pi_1(\overline X_y,\bar{a})\to \pi_1(X,a)\to \pi_1(Y,b)\to e.
$$
\end{corollaire}

\begin{remarques}
\label{X.1.5}
On notera que la d\'emonstration de~\ref{X.1.3} fait intervenir de
fa\c con essentielle~\ref{X.1.2} et par l\`a le \og premier
th\'eor\`eme de comparaison\fg en g\'eom\'etrie
alg\'ebrico-formelle. Par contre, la th\'eorie de la descente de
l'expos\'e~\ref{IX}
n'est intervenue que par l'interm\'ediaire de IX~\ref{IX.3.4}, dont
une d\'emonstration directe dans le cas d'un morphisme \emph{propre}
$f\colon X\to Y$ tel que $f_*(\cal{O}_X)=\cal{O}_Y$ est facile. Soit
en effet $Y'$ \'etale sur~$Y$ et supposons que $X' = X\times_Y Y'$
soit somme disjointe de deux ouverts non vides, prouvons qu'il en est
de m\^eme de~$Y'$. En effet, on aura $Y'=\Spec(\cal{A})$, donc
$X'=\Spec(\cal{B})$ avec $\cal{B}=\cal{A}\otimes_{\cal{O}_Y}
\cal{O}_X$, et la d\'ecomposition de~$X'$ en somme directe
correspond \`a une d\'ecomposition de~$\cal{B}$ en produit de deux
Alg\`ebres non nulles $\cal{B}_1$ et $\cal{B}_2$. Comme
$f_*(\cal{O}_X)=\cal{O}_Y$ on conclut facilement
$f_*(\cal{B})=\cal{A}$, donc $\cal{A}$ sera somme de deux Alg\`ebres
(non nulles \'egalement, car leurs sections unit\'e sont non
nulles) $f_*(\cal{B}_1)$ et $f_*(\cal{B}_2)$, cqfd.
\end{remarques}

\subsection{}
\label{X.I.6}
Supposons encore que $f$ soit propre et s\'eparable, mais ne faisons
plus d'hypoth\`ese sur~$f_*(\cal{O}_X)$, qui correspondra \`a un
rev\^etement \'etale $Y'$ bien d\'etermin\'e de~$Y$,
d'ailleurs ponctu\'e au-dessus de~$b$ par l'image $b'$
de~$a$. Appliquant alors \ref{X.1.4} au morphisme canonique $X\to Y'$,
et supposant $f$ surjectif, la suite exacte~\ref{X.1.4} est
remplac\'ee par la suivante, analogue de la suite exacte d'homotopie
des espaces fibr\'es en topologies alg\'ebriques:
$$
\pi_1(\overline X_y,\overline a)\to \pi_1(X,a)\to \pi_1(Y,b)\to
\pi_0(\overline X_y,\overline a)\to \pi_0(X,a)\to \pi_0(Y,b)\to e.
$$
Bien entendu, dans~\ref{X.1.4} on ne peut pas en g\'en\'eral
affirmer que l'homomorphisme
$\pi_1(\overline X_y,\overline a)\to \pi_1(X,a)$
soit injectif; en topologie alg\'ebrique, son noyau est l'image de
$\pi_2(Y,b)$, et il y aurait lieu en g\'eom\'etrie alg\'ebrique
\'egalement d'introduire des groupes d'homotopie en toutes
dimensions, et la suite exacte d'homotopie compl\`ete pour un
morphisme propre satisfaisant des hypoth\`eses convenables (par
exemple d'\^etre un morphisme
lisse).
On ne dispose \`a l'heure
actuelle d'aucun r\'esultat dans ce sens, \`a l'exception d'une
d\'efinition raisonnable (sinon d\'efinitive) des groupes
d'homotopie sup\'erieure.

\begin{corollaire}
\label{X.1.7}
Soient $k$ un corps alg\'ebriquement clos, $X$ et $Y$ deux
pr\'esch\'emas connexes sur~$k$, on suppose $X$ propre sur~$k$ et
$Y$ localement
noeth\'erien. Soient $a$ un point g\'eom\'etrique de~$X$, $b$
un point g\'eom\'etrique de~$Y$ \`a valeurs dans la m\^eme
extension alg\'ebriquement close $K$ de~$k$, consid\'erons le
point g\'eom\'etrique $c=(a,b)$ de~$X\times_k Y$, et
l'homomorphisme $\pi_1(X\times_k Y,c)\to \pi_1(X,a)\times\pi_1(Y,b)$
d\'eduit des homomorphismes sur les groupes fondamentaux
associ\'es aux deux projections $X\times_k Y\to X$ et $X\times_k
Y\to Y$. L'homomorphisme pr\'ec\'edent est un \emph{isomorphisme}.
\end{corollaire}

Supposons d'abord $K=k$. Posons $Z=X\times_k Y$, consid\'erons la
projection $f\colon Z\to Y$ et la localit\'e $y$ du point
g\'eom\'etrique $b$ de~$Y$, appliquons \`a la situation le
r\'esultat~\ref{X.1.4}. On notera pour ceci que quitte \`a passer
\`a $X_\red$ (ce qui ne change pas les groupes fondamentaux
envisag\'es), on peut supposer d\'ej\`a $X$ r\'eduit donc
s\'eparable sur~$k$, donc~$Z$ est s\'eparable sur~$k$, et
\'evidemment \`a fibres g\'eom\'etriquement connexes (puisque~$X$ est connexe). La fibre g\'eom\'etrique de~$Z$ en $b$ est
canoniquement isomorphe \`a $X\otimes_k K=X$. D'autre part, comme le
compos\'e des morphismes $X\to Z\to X$ est l'identit\'e, on trouve
que $\pi_1(X,a)\to \pi_1(Z,c)$ est injectif et 1.4 nous donne une
suite exacte:
$$
e\to \pi_1(X,a)\to \pi_1(Z,c)\to \pi_1(Y,b)\to e.
$$
D'autre part, on a la suite exacte canonique:
$$
e\to \pi_1(X,a)\to \pi_1(X,a)\times \pi_1(Y,b)\to \pi_1(Y,b)\to e,
$$
o\`u les deux homomorphismes \'ecrits sont l'injection canonique
et la projection canonique. On a enfin un homomorphisme de la
premi\`ere suite exacte dans la deuxi\`eme, \`a l'aide des
morphismes identiques sur les termes extr\^emes, et l'homomorphisme
canonique $\pi_1(Z,c)\to \pi_1(X,a)\times\pi_1(Y,b)$ pour les termes
m\'edians. La commutativit\'e du diagramme ainsi obtenu se
v\'erifie trivialement. Comme les homomorphismes sur les termes
extr\^emes sont des isomorphismes, il en est de m\^eme pour les
termes m\'edians, ce qui prouve \ref{X.1.7} dans ce cas.

Lorsqu'on ne suppose plus $K=k$, on trouve seulement un isomorphisme
$$
\pi_1(Z,c)\to \pi_1(X\otimes_k K,a)\times \pi_1(Y,b),
$$
et~\ref{X.1.7} \'equivaut alors au cas particulier suivant:

\begin{corollaire}
\label{X.1.8}
Soient $X$ un sch\'ema propre et connexe sur un corps
alg\'ebriquement clos~$k$, $k'$ une extension alg\'ebriquement
close de~$k$, $a'$ un point g\'eom\'etrique de~$X\otimes_k k'$ et
$a$ son image dans~$X$. Alors l'homomorphisme canonique
$\pi_1(X\otimes_k k',a')\to \pi_1(X,a)$ est un \emph{isomorphisme.}
\end{corollaire}

Le fait que cet homomorphisme soit surjectif \'equivaut \`a dire
que si $X'$ est un rev\^etement \'etale connexe de~$X$, alors
$X'\otimes_k k'$ est \'egalement connexe, et r\'esulte
aussit\^ot du fait que $k$ est alg\'ebriquement clos; c'est aussi
un cas particulier
de~IX~\ref{IX.3.4}.
L'hypoth\`ese de propret\'e sur~$X$ n'a pas encore servi. Ceci
dit, dire que l'homomorphisme envisag\'e est injectif signifie aussi
ceci: \emph{tout revêtement étale de~$X\otimes_k k'$ est
isomorphe à l'image inverse d'un revêtement étale de~$X$.}
Il est essentiellement sorital qu'on peut trouver une
sous-$k$-alg\`ebre $A$ de
$k'$,
de type fini sur~$k$, et un
rev\^etement \'etale de~$X\otimes_k A$ dont l'image inverse sur~$X\otimes_k k'$ soit isomorphe au rev\^etement donn\'e. Soit donc
$Y=\Spec(A)$, qui est un $k$-sch\'ema int\`egre de type fini, donc
a des points rationnels sur~$k$. Appliquons alors \ref{X.1.7} au
groupe fondamental de~$X\times Y$ en un point $(a,b)$ rationnel
sur~$k$: on trouve que tout rev\^etement \'etale connexe de
$X\times Y$ est isomorphe \`a un quotient d'un rev\^etement
$X'\times Y'$, o\`u $X'$, $Y'$ sont des rev\^etements galoisiens
\'etales de~$X$ et $Y$ de groupes~$G$, $G'$ par un sous-groupe $H$
de~$G\times G'$. Cela implique que l'image inverse de ce
rev\^etement de~$X\times Y$ sur~$X\times Y'$ est isomorphe \`a un
rev\^etement de la forme $X_1'\times Y'$, o\`u $X_1'$ est un
rev\^etement \'etale de~$X$. Si donc $L$ est le corps des
fonctions de~$Y$, \'egal au corps des fractions de~$A$ dans~$k'$, le
rev\^etement \'etale de~$X\otimes_k L$ induit par le
rev\^etement donn\'e de~$X\times_k Y$ est tel qu'il existe une
extension finie s\'eparable $L'$ de~$L$, telle que l'image inverse
dudit rev\^etement sur~$X\otimes_k L'$ est isomorphe \`a
$X_1'\otimes_k L'$. Or $k'$ \'etant alg\'ebriquement clos, on peut
supposer que l'extension $L'$ de~$L$ est contenue dans~$k'$. Cela
prouve que le rev\^etement \'etale donn\'e de~$X\otimes_k k'$
est isomorphe \`a $X_1'\otimes_k k'$, cqfd.

La forme explicite signal\'ee en passant pour les rev\^etements
\'etales d'un produit
$X\times_k Y$ implique aussit\^ot le r\'esultat suivant:

\begin{corollaire}
\label{X.1.9}
Soient $k$ un corps alg\'ebriquement clos, $X$ et $Y$ deux
pr\'esch\'emas localement noeth\'eriens sur~$k$, $Z=X\times_k
Y$ leur produit, $Z'$ un rev\^etement \'etale de~$Z$. Pour tout
point $y\in Y$ rationnel sur~$k$, soit $i_y \colon \Spec(k)\to Y$ le
morphisme canoniquement associ\'e, $j_y = \id_X\times_k i_y$ le
morphisme $X\to Z$ correspondant. Soit enfin $X_y'$ le rev\^etement
\'etale de~$X$ image inverse de~$Z'$ par~$j_y$. On suppose $Y$
connexe, et $X$ ou~$Y$ propre sur~$k$. Alors les rev\^etements
$X_y'$ de~$X$ sont tous isomorphes.
\end{corollaire}

De fa\c con imag\'ee, on peut dire qu'\emph{une famille de
revêtements étales de~$X$, paramétrée par un
préschéma connexe~$Y$, est constante si $X$ ou le
préschéma de paramètres~$Y$ est propre sur~$k$.}

\begin{remarques}
\label{X.1.10}
Les corollaires~\ref{X.1.7} \`a~\ref{X.1.9} sont d\^us \`a
Lang-Serre~\cite{X.2} dans le cas des sch\'emas alg\'ebriques normaux.
(Leur travail a \'et\'e la motivation initiale pour la th\'eorie
du groupe fondamental d\'evelopp\'ee dans ce S\'eminaire). Comme
l'ont remarqu\'e ces auteurs, ces r\'esultats deviennent inexacts
lorsqu'on y abandonne l'hypoth\`ese de propret\'e, du moins en
caract\'eristique $p>0$. Prenant par exemple pour $X$ la droite
affine $X=\Spec(k[t])$, il n'est pas difficile de voir que les
rev\^etements de~$X$, param\'etr\'es par la droite affine
$Y=\Spec(k[s])$, d\'efinis par les \'equations
$$
x^p - x = st,
$$
sont \'etales et deux \`a deux non isomorphes. Cela met en
d\'efaut~\ref{X.1.9} et a fortiori~\ref{X.1.7}, et on voit de
m\^eme que si $s$ est consid\'er\'e comme un \'el\'ement
transcendant sur~$k$ dans une extension alg\'ebriquement close $K$
de~$k$, on trouve un rev\^etement \'etale $X'$ de~$X\otimes_kK$
qui ne provient pas d'un rev\^etement \'etale de~$X$.
\end{remarques}

\section[Application du th\'eor\`eme d'existence de faisceaux]{Application du th\'eor\`eme d'existence de faisceaux:
th\'eo\-r\`eme de semi-con\-ti\-nu\-i\-t\'e pour les groupes
fondamentaux des fibres d'un morphisme propre et s\'eparable}
\label{X.2}

\begin{theoreme}
\label{X.2.1}
Soient $Y$ le spectre d'un anneau local noeth\'erien
\emph{complet}, de corps r\'esiduel~$\kres$, $X$ un $Y$-sch\'ema
propre, $X_0=X\otimes_A \kres$, $a_0$ un point g\'eom\'etrique de
$X_0$ et $a$ le point g\'eom\'etrique correspondant de~$X$. Alors
l'homomorphisme canonique $\pi_1(X_0,a_0)\to \pi_1(X,a)$ est un
\emph{isomorphisme.}
\end{theoreme}

Ce n'est qu'une traduction, dans le langage du groupe fondamental, du
r\'esultat rappel\'e dans IX~\ref{IX.1.10}. C'est ici que le
th\'eor\`eme d'existence
de
faisceaux en g\'eom\'etrie
alg\'ebrico-formelle s'introduit de fa\c con essentielle dans la
th\'eorie du groupe fondamental.

Introduisons maintenant une cl\^oture alg\'ebrique $\bar{k}$ du
corps r\'esiduel $k$, et la fibre g\'eom\'etrique $\overline X_0
= X_0\otimes_k \bar{k}$. On a donc la suite exacte (IX~\ref{IX.6.1})
$$
e\to \pi_1(\overline X_0,\overline a)\to \pi_1(X_0,a_0)\to
\pi_1(k,\bar{k})\to e,
$$
d'autre part on a l'isomorphisme~\ref{X.2.1} et l'isomorphisme
analogue, plus \'el\'ementaire,
$$
\pi_1(k,\bar{k})\to \pi_1(Y,b),
$$
o\`u $b$ est l'image de~$a$ dans~$Y$. On trouve ainsi:

\begin{corollaire}
\label{X.2.2}
Avec les notations pr\'ec\'edentes, supposons
$\overline X_0=X_0\otimes_k \bar{k}$
connexe, et soient $\overline a_0$ un point g\'eom\'etrique de
$\overline X_0$, $a_0$ son image dans~$X$,
$b_0$ son image dans~$Y$. Alors la suite d'homomorphismes canoniques
suivante
$$
e\to \pi_1(\overline X_0,\overline a)\to \pi_1(X,a_0)\to
\pi_1(Y,b_0)\to e
$$
est exacte.
\end{corollaire}

On comparera cette suite \`a la suite exacte~\ref{X.1.4}, mais on
notera que~a)~on n'a pas eu \`a faire d'hypoth\`ese de platitude,
ou de s\'eparabilit\'e sur les fibres, pour $X\to Y$; b)~on a le
compl\'ement important que \emph{le morphisme $\pi_1(\overline
X_0,\overline a_0)\to \pi_1(X,a_0)$ est injectif.}

Ce dernier fait nous permettra de comparer le groupe fondamental des
autres
fibres g\'eom\'etriques de~$X$ sur~$Y$ \`a celui de~$\overline
X_0$. Soit en effet $y_1$ un point quelconque de~$Y$, $X_1$ sa fibre
et $\overline X_1$ sa fibre g\'eom\'etrique, relativement \`a
une extension alg\'ebriquement close de~$\kres(y_1)$, $\overline a_1$
un point g\'eom\'etrique de~$\overline X_1$, $a_1$ son image dans
$X$ et $b_1$ son image dans~$Y$. Choisissons une \og classe de chemins\fg
de~$a_1$ \`a $a_0$, d'o\`u une classe de chemins de~$b_1$ \`a
$b_0$, d'o\`u un diagramme commutatif d'homomorphismes:
$$
\xymatrix@=.5cm{
&{\pi_1(\overline X_1,\overline a_1)}\ar[r]&{\pi_1(X,a_1)}\ar[r]\ar[d]&{\pi_1(Y,b_1)}\ar[r]\ar[d]&{e}\\{e}\ar[r]&{\pi_1(\overline X_0,\overline a_0)}\ar[r]&{\pi_1(X,a_0)}\ar[r]&{\pi_1(Y,b_0)}\ar[r]&{e,}}
$$
o\`u les deux fl\`eches verticales \'ecrites sont des
isomorphismes. Comme la deuxi\`eme ligne est exacte, on trouve donc
un homomorphisme canonique, que nous appellerons \emph{l'homomorphisme
de spécialisation pour le groupe fondamental} (ne d\'ependant
que de la classe de chemins choisis de~$a_1$ \`a~$a_0$, donc
\emph{défini modulo automorphisme intérieur de}
$\pi_1(X,a_0)$):
$$
\pi_1(\overline X_1,\overline a_1)\to \pi_1(\overline X_0,\overline
a_0).
$$
Lorsque la premi\`ere ligne ci-dessus est \'egalement exacte, il
s'ensuit aussit\^ot que l'homo\-mor\-phisme de sp\'ecialisation
est surjectif. On trouve donc, compte tenu de~\ref{X.1.4}:

\begin{corollaire}
\label{X.2.3}
Sous les conditions de~\ref{X.2.1}, supposons de plus que le morphisme
$f\colon X\to Y$ soit \emph{séparable} \eqref{X.1.1} et $\overline
X_0$ connexe (donc en vertu de~\ref{X.1.2} on a
$f_*(\cal{O}_X)=\cal{O}_Y$). Alors pour toute fibre
g\'eom\'etrique $\overline X_1$ de~$X$ sur~$Y$, munie d'un point
g\'eom\'etrique $\overline a_1$, l'homomorphisme de
sp\'ecialisation
\index{homomorphisme de sp\'ecialisation pour le groupe fondamental|hyperpage}\index{sp\'ecialisation pour le groupe fondamental (homomorphisme de)|hyperpage}d\'efini ci-dessus est un homomorphisme \emph{surjectif.}
\end{corollaire}
C'est l\`a un r\'esultat de \emph{semi-continuité} pour le
groupe fondamental, qui ne semble pas encore avoir d'analogue en
topologie alg\'ebrique. On peut d'ailleurs l'\'enoncer sous des
conditions plus g\'en\'erales:

\begin{corollaire}
\label{X.2.4}
Soient
$f\colon X\to Y$ un morphisme propre \`a fibres
g\'eom\'etriquement connexes, avec $Y$ localement noeth\'erien,
$y_0$ et $y_1$ deux points de~$Y$ tels que $y_0\in\overline{\{y_1\}}$,
$\overline X_0$ et $\overline X_1$ les fibres g\'eom\'etrique de
$X$ correspondant \`a des extensions alg\'ebriquement closes
donn\'ees de~$\kres(y_0)$ et $\kres(y_1)$, $\overline a_0$
\resp $\overline a_1$ un point g\'eom\'etrique de~$\overline X_0$
\resp $\overline X_1$. Alors on peut d\'efinir de fa\c con
naturelle un homomorphisme de sp\'ecialisation:
$$
\pi_1(\overline X_1,\overline a_1)\to \pi_1(\overline X_0,\overline
a_0),
$$
d\'efini \`a automorphisme int\'erieur pr\`es, et c'est l\`a
un homomorphisme \emph{surjectif} si $f$ est un morphisme
s\'eparable \eqref{X.1.1}.
\end{corollaire}

En effet, il r\'esulte d'abord de~\ref{X.1.8} que \ref{X.2.4} est
essentiellement ind\'ependant des extensions alg\'ebriquement
closes choisies pour les corps r\'esiduels $\kres(y_0)$
et~$\kres(y_1)$. Cela nous permet de remplacer $Y$ par un
pr\'esch\'ema $Y'$ sur~$Y$ ayant un point $y_0'$ (\resp~$y_1'$)
au-dessus de~$y_0$ (\resp $y_1$). On prendra alors pour $Y'$ le
spectre du compl\'et\'e de l'anneau local de~$y_0$ dans~$Y$, et on
applique~\ref{X.2.3}.

\begin{remarques}
\label{X.2.5}
La conclusion finale de~\ref{X.2.4} sur la surjectivit\'e de
l'homomorphisme de sp\'e\-cia\-li\-sa\-tion, et a fortiori les
r\'esultats~\ref{X.1.3} et~\ref{X.1.4} dont elle est une
cons\'equence, devient inexacte si on ne suppose plus que $f\colon
X\to Y$ est un morphisme s\'eparable, m\^eme pour des sch\'emas
projectifs sur un corps alg\'ebriquement clos de
caract\'eristique~$0$. Nous en verrons plus loin des exemples, tant
dans le cas o\`u $f$ est plat mais o\`u $f$ admet une fibre non
s\'eparable ($X$ et $Y$ \'etant cependant lisses sur~$k$), que
dans le cas o\`u les fibres de~$f$ sont bien s\'eparables mais
o\`u $f$ n'est pas plat (par exemple $f\colon X\to Y$ \'etant un
morphisme birationnel de sch\'emas int\`egres normaux),
\cf XI~\ref{XI.3}. Dans ces exemples, il peut arriver que le groupe
fondamental de la fibre g\'eom\'etrique g\'en\'erique soit
nul, mais non celui d'une fibre g\'eom\'etrique sp\'eciale
convenable. D'autre part, m\^eme si $f\colon X\to Y$ est un
morphisme propre s\'eparable comme dans~\ref{X.2.4}, il arrive
couramment que le morphisme de sp\'ecialisation ne soit
pas un isomorphisme. Ainsi, il est facile de donner des exemples
o\`u $\overline X_1$ est une courbe elliptique non singuli\`ere,
(donc son groupe fondamental est commutatif, et sa composante
$\ell$-primaire pour un nombre premier $\ell$ premier \`a la
caract\'eristique est isomorphe \`a $\ZZ_\ell^2$, \cf \ref{XI}),
tandis que $\overline X_0$ est form\'e, soit de deux courbes
rationnelles non singuli\`eres se coupant en deux points, soit de
deux courbes rationnelles tangentes en un point, soit enfin d'une
courbe rationnelle ayant un point singulier qui est un point de
rebroussement. (Pour la classification compl\`ete des courbes
elliptiques d\'eg\'en\'er\'ees, voir les travaux r\'ecents
de Kodaira~\cite{X.1} et N\'eron). On voit alors que dans ces cas, le
groupe fondamental de~$\overline X_0$ est respectivement $\hat\ZZ$,
$e$, $e$, donc \og strictement plus petit\fg que celui de~$\overline
X_1$. Nous verrons cependant plus loin, lorsque $f$ est un morphisme
\emph{lisse}, une majoration du noyau de l'homomorphisme de
sp\'ecialisation, qui implique en particulier que si $\kres(y_0)$ est
de \emph{caractéristique} 0, l'homomorphisme de sp\'ecialisation
est un isomorphisme. Mais m\^eme pour un morphisme lisse, si la
caract\'eristique de~$\kres(y_0)$ est $>0$, il peut arriver que
l'homomorphisme de sp\'ecialisation ne soit pas un isomorphisme,
comme on voit par exemple dans le cas o\`u $X$ est un sch\'ema
ab\'elien sur~$Y$ (de dimension relative~$1$, si on veut),
\cf XI~\ref{XI.2}.
Une th\'eorie satisfaisante de la
sp\'ecialisation du groupe fondamental doit tenir compte de la
\og composante continue\fg du \og vrai\fg groupe fondamental, correspondant
\`a la classification des rev\^etements principaux de groupe
structural des groupes infinit\'esimaux; moyennant quoi on serait en
droit \`a s'attendre que les \og vrais\fg groupes fondamentaux des
fibres g\'eom\'etriques d'un morphisme lisse et propre $f\colon
X\to Y$ forment un joli syst\`eme local sur~$X$, limite projective
de sch\'emas en groupes finis et plats sur~$X$\footnote{Cette
conjecture extr\^emement s\'eduisante est malheureusement mise en
d\'efaut par un exemple in\'edit de M.\, Artin, d\'ej\`a lorsque
les fibres de~$f$ sont des courbes alg\'ebriques de genre donn\'e
$g\ge 2$.}. Nous reviendrons ult\'erieurement sur ce point de vue,
notre objet pr\'esent \'etant au contraire de pousser aussi loin
que possible les ph\'enom\`enes communs \`a la th\'eorie
topologique et la th\'eorie sch\'ematique du groupe fondamental.
\end{remarques}

Soit maintenant $X_0$ une courbe propre, lisse et connexe de genre $g$
sur un corps al\-g\'e\-bri\-quement clos~$k$. Si $k$ est de
caract\'eristique z\'ero, son groupe fondamental peut se
d\'eterminer par voie transcendante de la fa\c con suivante. On
sait que $X_0$ provient par extension de la base d'une courbe
d\'efinie
sur une extension alg\'ebriquement close de degr\'e de
transcendance fini du corps premier $\QQ$, et compte tenu
de~\ref{X.1.8}, on peut supposer que $k$ est lui-m\^eme de degr\'e
de transcendance fini sur~$\QQ$. On peut donc supposer que $k$ est un
sous-corps du corps $\CC$ des nombres complexes, et une nouvelle
application de~\ref{X.1.8} nous permet de supposer que $k=\CC$. Il
n'est pas difficile alors de v\'erifier que le groupe fondamental de
$X$ est isomorphe au compactifi\'e du groupe fondamental de l'espace
topologique associ\'e $\tilde X$ (surface compacte orient\'ee de
genre~$g$), pour la topologie d\'efinie par les sous-groupes
d'indice fini\footnote{Cette d\'eduction \'etait explicit\'ee
dans un des expos\'es oraux qui n'ont pas \'et\'e
r\'edig\'es.}. Il est d'autre part classique que le groupe
fondamental topologique est engendr\'e par $2g$ g\'en\'erateurs
$s_i,t_i$ ($1\le i\le g$), soumis \`a une seule relation:
$$
(s_1t_1s_1^{-1}t_1^{-1})\cdots(s_gt_gs_g^{-1}t_g^{-1})=1.
$$
Donc le groupe fondamental de~$X$ admet $2g$ g\'en\'erateurs
\emph{topologiques} $s_i,t_i$ ($1\le i\le g$), li\'es par la seule
relation pr\'ec\'edente. Si maintenant la caract\'eristique de
$k$ est $p>0$, d\'esignons par $A$ l'anneau des vecteurs de Witt
construit avec~$k$, par $K$ une extension alg\'ebriquement close de
son corps des fractions. On a vu dans III~\ref{III.7.4} qu'il
existe un sch\'ema $X$ sur~$Y=\Spec(A)$, propre et lisse sur~$Y$, se
r\'eduisant suivant $X_0$. Appliquons-lui~\ref{X.2.3}, on trouve un
morphisme \emph{surjectif}
$$
\pi_1(X_1)\to \pi_1(X_0),
$$
o\`u $X_1=X\otimes_A K$. Il est imm\'ediat\footnote{\cf EGA
IV~12.2.} que $X_1$ est lisse sur~$K$, connexe \eqref{X.1.2}, de
dimension~$1$, et son genre est \'egal \`a~$g$ (d'apr\`es
l'invariance de la caract\'eristique d'Euler-Poincar\'e, \cf EGA
III~7). Comme $K$ est de caract\'eristique~$0$, on peut lui
appliquer le r\'esultat pr\'ec\'edent. On a ainsi prouv\'e par
\emph{voie transcendante}:

\begin{theoreme}
\label{X.2.6}
Soit $X_0$ une courbe alg\'ebrique lisse propre et connexe sur un
corps al\-g\'e\-bri\-quement clos $k$, et soit $g$ son genre. Alors
$\pi_1(X_0)$ admet un syst\`eme de~$2g$ g\'en\'erateurs
topologiques, li\'es par la relation \'ecrite plus haut. Lorsque la
caract\'eristique de~$k$ est~$0$, $\pi_1(X_0)$ est m\^eme le
groupe de type galoisien libre pour les g\'en\'erateurs et la
relation qui pr\'ec\`edent.
\end{theoreme}

\begin{remarques}
\label{X.2.7}
Il n'existe pas \`a l'heure actuelle, \`a la connaissance du
r\'edacteur, de d\'e\-mon\-stra\-tion par voie purement
alg\'ebrique du r\'esultat pr\'ec\'edent, (sauf pour les genres
$0,1$). Pour commencer, on ne voit gu\`ere comment distinguer dans
$\pi_1(X_0)$ $2g$~\'el\'ements, dont on pourrait attendre ensuite
qu'ils forment un syst\`eme de g\'en\'erateurs
topologiques. \`A cet \'egard, la situation de la droite
rationnelle priv\'ee de~$n$ points, et l'\'etude des
rev\^etements d'icelle mod\'er\'ement ramifi\'es en ces
points, est plus sympathique, puisque la consid\'eration des groupes
de ramification en ces $n$ points fournit $n$ \'el\'ements du
groupe fondamental \`a \'etudier, dont on montre en effet qu'ils
engendrent topologiquement ce groupe fondamental, comme nous verrons
ult\'erieurement\footnote{\Cf Exp.\, \ref{XII}. Encore ces
\'el\'ements ne sont-ils d\'etermin\'es vraiment que modulo
conjugaison, et il convient de faire un choix \emph{simultané}
judicieux de ces \'el\'ements dans leurs classes.}. Mais m\^eme
dans ce cas particuli\`erement concret, il ne semble pas exister de
d\'emonstration purement alg\'ebrique. Une telle
d\'emonstration serait \'evidemment extr\^emement
int\'eressante. Le seul fait concernant le groupe fondamental d'une
courbe qu'on sache d\'emontrer par voie purement alg\'ebrique
(exception faite du th\'eor\`eme de finitude faible~\ref{X.2.12}
ci-dessous, prouv\'e par voie alg\'ebrique par Lang-Serre~\cite{X.2}),
semble la d\'etermination du groupe fondamental rendu ab\'elien
via la jacobienne (signal\'ee dans IX~\ref{IX.5.8} derni\`ere
ligne).
\end{remarques}

\subsection{}
\label{X.2.8}
La derni\`ere assertion~\ref{X.2.6} n'est plus valable en
caract\'eristique $p>0$, comme on voit d\'ej\`a dans le cas des
courbes elliptiques. Comme nous l'avons d\'ej\`a signal\'e, nous
ne savons pas si le groupe fondamental de~$X_0$ est topologiquement de
pr\'esentation finie; cela semble tout \`a fait improbable.

\begin{theoreme}
\label{X.2.9}
Soient $k$ un corps alg\'ebriquement clos, et $X$ un sch\'ema
propre et connexe sur~$k$. Alors le groupe fondamental de~$X$ est
topologiquement de g\'en\'eration finie.
\end{theoreme}

Nous
proc\'edons par r\'ecurrence sur~$n=\dim X$, l'assertion
\'etant triviale pour~$n\le 0$. Supposons donc~$n>0$, et le
th\'eor\`eme d\'emontr\'e pour les
dimensions~$n'<n$. D'apr\`es le lemme de Chow (EGA~II~5.6.2) il
existe un sch\'ema projectif $X'$ sur~$k$ et un morphisme surjectif
$X'\to X$. On peut \'evidemment supposer $X'$ r\'eduit, et en
passant au normalis\'e, normal. Gr\^ace \`a la th\'eorie de
la descente, il suffit de prouver que les groupes fondamentaux des
composantes connexes de~$X'$ sont topologiquement de
g\'en\'eration finie (IX~\ref{IX.5.2}). Cela nous ram\`ene
donc au cas o\`u $X$ est \emph{projectif} et \emph{normal}. Si
alors $n=1$, il suffit d'appliquer~\ref{X.2.6}. Si $n\ge 2$, on
consid\`ere une immersion projective $X\to \PP_k^r$, et une section
hyperplane $Y=X\cdot H$
(munie
de la structure r\'eduite induite), telle que
$Y\ne X$,
\ie $H\not\supset X$. On aura alors $\dim Y<n$, et tenant compte de
l'hypoth\`ese de r\'ecurrence, il suffit de prouver que
$\pi_1(Y)\to \pi_1(X)$ est \emph{surjectif.} Or plus
g\'en\'eralement:

\begin{lemme}
\label{X.2.10}
Soient $X$ un pr\'esch\'ema propre sur un corps alg\'ebriquement
clos~$k$, $g\colon X\to \PP_k^r$ un morphisme. On suppose $X$
irr\'eductible et normal et $\dim g(X)\ge 2$. Soient $H$ un
hyperplan de~$\PP_k^r$ et $Y=X\times_{\PP_k^r} H$. Alors $Y$ est
connexe, et l'homomorphisme $\pi_1(Y)\to\pi_1(X)$ est surjectif.
\end{lemme}

Ces assertions r\'esultent en effet de la suivante:

\begin{corollaire}
\label{X.2.11}
Sous les conditions pr\'ec\'edentes, soient $X'$ un rev\^etement
\'etale connexe de~$X$, et $Y'=X'\times_X Y=X'\times_{\PP_k^r} H$ le
rev\^etement induit sur~$Y$. Alors $Y'$ est connexe.
\end{corollaire}

Comme $X$ est normal, $X'$ est normal, donc \'etant connexe, $X'$
est irr\'eductible; de plus son image dans $\PP_k^r$ est de
dimension~$\ge 2$. Un lemme bien connu d\^u \`a Zariski (et
appel\'e \og \emph{théorème de Bertini}\fg) implique donc que
si $H_1'$ est l'hyperplan g\'en\'erique dans~$\PP_k^r$,
d\'efini sur une extension $K$ de~$k$, alors $X'\times_{\PP^r} H_1$
est universellement irr\'eductible donc universellement connexe
sur~$K$. Le th\'eor\`eme de connexion de Zariski (EGA~III~4)
implique alors que pour \emph{tout} hyperplan $H$ (d\'efini sur une
extension quelconque de~$k$) $X'\times_{\PP^r} H$ est
g\'eom\'etriquement connexe. Cela ach\`eve la d\'emonstration
de~\ref{X.2.11}, donc de~\ref{X.2.9}.

\begin{corollaire}[Lang-Serre]
\label{X.2.12}
Sous
les conditions de~\ref{X.2.9}, pour tout
groupe fini~$G$, l'ensemble des classes, \`a isomorphisme pr\`es,
de rev\^etements principaux de~$X$ de groupe~$G$, est fini.
\end{corollaire}

\begin{remarque}
\label{X.2.13}
Sous les conditions de~\ref{X.2.10} nous prouverons lorsque $\dim
g(X)\ge 3$ (du moins lorsque $g$ est une immersion et $X$
r\'egulier), un r\'esultat plus pr\'ecis, connu en
g\'eom\'etrie alg\'ebrique sous le nom de
\og \emph{théorème de Lefschetz}\fg: $\pi_1(Y)\to \pi_1(X)$
\emph{est un isomorphisme.}\footnote{\Cf le s\'eminaire SGA~2 (1962)
faisant suite \`a celui-ci, th\'eor\`eme X~3.10.}
Il y a dans les cas classiques des
\'enonc\'es analogues pour les groupes d'homologie et les groupes
d'homotopie sup\'erieure, qui t\^ot ou tard devront \^etre
englob\'es dans la g\'eom\'etrie alg\'ebrique
abstraite. M\^eme pour la cohomologie de Hodge $\H^p(X,\Omega^q)$,
il ne semble pas que la question ait encore \'et\'e
\'etudi\'ee; il n'est d'ailleurs gu\`ere probable que pour cette
derni\`ere, les th\'eor\`emes de Lefschetz subsistent tels quels
en caract\'eristique~$p>0$.
\end{remarque}

\begin{remarqueMR}
\label{X.2.14}
Soit $R$ un anneau de valuation discr\`ete complet, de
corps r\'esiduel $\kres$ alg\'ebriquement clos de caract\'eristique $p>0$, de
corps des fractions~$K$ et soit $Y$ une courbe propre et lisse,
connexe, de genre $g$ sur~$R$. On a donc un morphisme de
sp\'ecialisation surjectif $\sp : \pi_1(Y_{\overline{K}}) 
\to\pi_1(Y_{\kres})$.
On a d\'ej\`a not\'e que si $K$ est de caract\'eristique 0, $\sp$ n'est pas un
isomorphisme d\`es que $g\geq1$. Supposons~$K$ de caract\'eristique $p$,
de sorte que $R$ est isomorphe \`a l'anneau de s\'eries 
formelles~$\kres[[T]]$.
En \'egale caract\'eristique $p >0$, le groupe fondamental n'est pas
d\'etermin\'e par le genre~$g$ comme on le voit d\'ej\`a sur les courbes
elliptiques qui peuvent \^etre ordinaires ou supersinguli\`eres. Citons
le r\'esultat r\'ecent de A.\, Tamagawa (non encore publi\'e). Si~$G$ est un
groupe profini, on note $G^{\res}$, le groupe profini quotient de~$G$,
limite projective des quotients topologiques de~$G$ qui sont finis
r\'esolubles.

\begin{theoreme*}[A.\, Tamagawa]
Supposons $g\geq2$, que la fibre sp\'eciale $Y_{\kres}$
soit d\'efinissable sur un corps fini et que le morphisme
$\sp^{\res}:\pi_1(Y_{\overline{K}})^{\res}\to\pi_1(Y_{\kres})^{\res}$ d\'eduit de
$\sp$ par passage au quotient, soit bijectif. Alors la courbe $Y$ est
constante sur~$R$.
\end{theoreme*}

Notons que le groupe de Galois de~$\overline{K}/K$ est r\'esoluble. L'\'enonc\'e
pr\'ec\'edent peut encore se traduire de la fa\c{c}on suivante : supposons que
tout rev\^etement fini \'etale $X_K\to Y_K$ de la fibre g\'en\'erique $Y_K$,
galoisien de groupe de Galois r\'esoluble, ait potentiellement bonne
r\'eduction (c'est-\`a-dire s'\'etende en un rev\^etement fini \'etale de~$Y$
apr\`es remplacement \'eventuel de~$R$ par son normalis\'e dans une extension
finie de~$K$). Alors $Y$ est constante sur~$R$.
\end{remarqueMR}

\section[Application du th\'eor\`eme de puret\'e]{Application du th\'eor\`eme de puret\'e:
th\'eor\`eme de continuit\'e pour les groupes fondamentaux des
fibres d'un morphisme propre et lisse}
\label{X.3}

Rappelons sans d\'emonstration\footnote{Pour une d\'emonstration, \cf SGA~2 X~\ref{X.3.4}.} le

\begin{theoremedepurete}[Zariski-Nagata]
\label{X.3.1}
\index{theoreme de purete@th\'eor\`eme de puret\'e|hyperpage}Soit $f\colon X\to Y$ un morphisme
quasi-fini et dominant de pr\'esch\'emas int\`egres, avec $X$
normal, $Y$ r\'egulier localement noeth\'erien, et soit $Z$
l'ensemble des points de~$X$ o\`u $f$ n'est pas \'etale,
\ie o\`u~$f$ est ramifi\'e (cela revient au m\^eme,
\textup{I~\ref{I.9.5}~(ii)}). Si $Z\ne X$, $Z$ est de codimension~$1$ dans
$X$ en tous ses points, \ie pour toute composante irr\'eductible
$Z'$ de~$Z$ de point g\'en\'erique~$z$, la dimension de Krull de
$\cal{O}_{X,z}$ est \'egale \`a~$1$.
\end{theoremedepurete}

Rappelons qu'un pr\'esch\'ema est dit \emph{normal}
\resp \emph{régulier} si ses anneaux locaux sont normaux
\resp r\'eguliers, et que la relation $Z\ne X$ signifie aussi que
l'extension finie $R(Z)/R(X)$ (o\`u $R$ d\'esigne le corps des
fonctions rationnelles) est \emph{séparable}. Se pla\c cant en
le point g\'en\'erique $z$ d'une composante $Z'$ de~$Z$, et
localisant en le point $y$ de~$Y$ en-dessous de~$z$, on trouve
l'\'enonc\'e \'equivalent:

\begin{corollaire}
\label{X.3.2}
Soient $A$ un anneau local noeth\'erien r\'egulier, $A\to B$ un
homomorphisme local injectif tel que $B$ soit normal, localis\'e
d'une alg\`ebre de type fini sur~$A$, et \emph{quasi-fini} sur~$A$;
on suppose de plus que $\dim A (=\dim B)\ge 2$, et que pour tout
id\'eal premier $\goth{p}$ de~$B$ distinct de l'id\'eal maximal,
$B$ est \'etale sur~$A$ en~$\goth{p}$, \ie $B_\goth{p}$ est
\'etale sur~$A_\goth{q}$ (o\`u $\goth{q}=A\cap\goth{p}$). Alors
$B$ est \'etale sur~$A$.
\end{corollaire}

D'ailleurs, il n'est pas difficile de r\'eduire ce dernier
\'enonc\'e au cas o\`u $A$ est un anneau local \emph{complet},
donc o\`u $B$ est \emph{fini} sur~$A$. Zariski~\cite{X.5} donne une
d\'emonstration simple de ce r\'esultat, valable dans le cas
d'\'egales caract\'eristiques; le cas g\'en\'eral est d\^u
\`a Nagata~\cite{X.3}, qui s'appuie sur un r\'esultat d\'elicat de
Chow; ce dernier n'a \'et\'e v\'erifi\'e par aucun des
participants du S\'eminaire, et devrait faire l'objet d'un
expos\'e ult\'erieur. Signalons seulement ici la d\'emonstration
tr\`es simple dans le cas particulier o\`u $\dim A=2$, qui est
suffisant pour l'application la plus importante que nous en ferons
dans le pr\'esent num\'ero. Comme $B$ est normal, il est un
$B$-module de profondeur (ancienne terminologie: codimension
cohomologique) $\ge 2$, donc c'est un $A$-module de profondeur~$\ge
2$, et comme $A$ est r\'egulier de dimension $2$, il en r\'esulte
que $B$ est un \emph{module libre} sur~$A$\footnote{\Cf EGA $0_{\textup{IV}}$
17.3.4}. Il r\'esulte alors de I~\ref{I.4.10} que l'ensemble
des id\'eaux premiers~$\goth{q}$ de~$A$ en lesquels $B$ est
ramifi\'e sur~$A$ est la partie de~$\Spec(A)$ d\'efinie par un
id\'eal principal (engendr\'e par le discriminant d'une base de
$B$ sur~$A$), donc est vide si elle est contenue dans le point
ferm\'e de~$\Spec(A)$, ce qui prouve~\ref{X.3.2} lorsque $\dim A=2$.

Nous utiliserons surtout~\ref{X.3.1} sous la forme \'equivalente:

\begin{corollaire}
\label{X.3.3}
Soient $X$ un pr\'esch\'ema localement noeth\'erien
r\'egulier, $U$ une partie ouverte de~$X$ com\-pl\'e\-men\-taire
d'une partie ferm\'ee $Z$ de~$X$ de codimension~$\ge 2$. Alors le
foncteur $X'\mto X'\times_X U$ de la cat\'egorie des
rev\^etements \'etales de~$X$ dans la cat\'egorie des
rev\^etements \'etales de~$U$ est une \'equivalence
de cat\'egories; en particulier, si $a$ est un point
g\'eom\'etrique de~$U$, l'homomorphisme canonique $\pi_1(U,a)\to
\pi_1(X,a)$ est un isomorphisme.
\end{corollaire}

La derni\`ere assertion est \'evidemment cons\'equence de la
premi\`ere, et pour celle-ci on peut \'evidemment supposer que $X$
est connexe donc irr\'eductible. De la normalit\'e de~$X$
r\'esulte d\'ej\`a que le foncteur $X'\mto X'\times_X U$ de
la cat\'egorie des rev\^etements localement libres (pas
n\'ecessairement \'etales) de~$X'$ dans la cat\'egorie des
rev\^etements de~$U$ est pleinement fid\`ele, car le foncteur
$\cal{E}\mto\cal{E}\vert U$ de la cat\'egorie des Modules
localement libres sur~$X$ dans la cat\'egorie des Modules localement
libres sur~$U$ l'est. Il reste donc \`a prouver que pour tout
rev\^etement \emph{étale} $U'$ de~$U$, il existe un
rev\^etement \'etale $X'$ de~$X$ (n\'ecessairement unique
d'apr\`es ce qui pr\'ec\`ede), tel que $U'$ soit isomorphe \`a
$X'\times_X U$. On peut \'evidemment supposer $U'$ connexe, donc
irr\'eductible puisque ($U$~\'etant normal) $U'$ est normal. Soit
$K$ le corps des fonctions rationnelles sur~$X$, ou sur~$U$ (c'est
pareil), $K'$ celui de~$U'$, alors $U'$ s'identifie au normalis\'e
de~$U$ dans~$K'$ I~\ref{I.10.3}. Soit $X'$ le normalis\'e de
$X$ dans~$K'$ (EGA~II~6.3), alors $X'=X\times_X U \cong U'$, d'autre
part $X'$ est normal, int\`egre, le morphisme structural $f\colon
X'\to X$ est \emph{fini} et dominant (car $X$ est normal et $K'/K$ est
une extension finie s\'eparable). Il est \'etale dans
$U'=f^{-1}(U)=X'=f^{-1}(Z)$, et comme $Z$ est de codimension $\ge 2$
dans $X$, $f^{-1}(Z)$ est de codimension $\ge 2$ dans~$X'$. On conclut
alors de~\ref{X.3.1} que $X'$ est \'etale sur~$X$, ce qui ach\`eve
la d\'emonstration.

\medskip
Soit maintenant $f\colon X\to Y$ une application rationnelle d'un
pr\'esch\'ema localement noeth\'erien et r\'egulier $X$
dans un pr\'esch\'ema~$Y$, et supposons que $f$ soit d\'efini
dans un ouvert $U$ compl\'ementaire d'une partie ferm\'ee de
codimension~$\ge 2$. Alors on d\'eduit de~\ref{X.3.3} un foncteur,
d\'efini \`a isomorphisme pr\`es, de la cat\'egorie des
rev\^etements \'etales de~$Y$ dans la cat\'egorie des
rev\^etements \'etales de~$X$, d'o\`u pour tout point
g\'eom\'etrique $a$ de~$U$, d'image $b$ dans~$Y$, un homomorphisme
canonique
$$
\pi_1(X,a)\to \pi_1(Y,b)
$$
(d\'eduit
de l'homomorphisme canonique $\pi_1(U,a)\to \pi_1(Y,b)$ gr\^ace
\`a l'isomorphisme $\pi_1(U,a)\isomto \pi_1(X,a))$. Lorsque $f$ est
un morphisme dominant, $X$ et $Y$ \'etant int\`egres de corps $K$
et~$L$, de sorte que $K$ est une extension de~$L$, et que $Y$ est
normal, ces correspondances se pr\'ecisent en termes d'extensions de
corps en notant que pour toute extension finie $L'$ de~$L$, non
ramifi\'ee sur~$Y$, l'alg\`ebre $K'=L'\otimes_L K$ sur~$K$ est non
ramifi\'ee sur~$X$.

En particulier, ces r\'eflexions montrent que le groupe fondamental
des pr\'esch\'emas localement noeth\'eriens connexes
r\'eguliers, ponctu\'es par des points g\'eom\'etriques
localis\'es en codimension~$\le 1$, est un \emph{foncteur} lorsque
l'on prend comme morphismes dans cette cat\'egorie les applications
rationnelles dominantes d\'efinies dans des compl\'ementaires de
parties ferm\'ees de codimension~$\ge 2$. Se rappelant par exemple
qu'une application rationnelle d'un sch\'ema normal sur un corps $k$
dans un sch\'ema propre sur~$k$ est d\'efinie dans le
compl\'ementaire d'un ensemble de codimension~$\ge 2$, on trouve:

\begin{corollaire}[Invariance birationnelle du groupe fondamental]
\label{X.3.4}
\index{invariance birationnelle du groupe fondamental|hyperpage}\index{theoreme d'invariance birationnelle du groupe fondamental@th\'eor\`eme d'invariance birationnelle du groupe fondamental|hyperpage}Soient $k$ un corps, $X$ et $Y$ deux sch\'emas propres sur~$k$ et
r\'eguliers, $f\colon X\to Y$ une application birationnelle de~$X$
dans~$Y$, $\Omega$ une extension alg\'ebriquement close du corps des
fonctions $K$ de~$X$, permettant de d\'efinir le groupe fondamental
de~$X$ et le groupe fondamental de~$Y$. Ces derniers sont alors
canoniquement isomorphes.
\end{corollaire}

Cela signifie aussi que pour une extension finie $K'$ de~$K$, si elle
est non ramifi\'ee sur un \og mod\`ele\fg propre non singulier $X$
de~$K$, elle l'est sur tout autre mod\`ele propre non singulier.

\begin{remarque}
\label{X.3.5}
Pour d'autre applications du th\'eor\`eme de puret\'e, voir les
travaux
d'\textsc{Abhyankar}
expos\'es dans~\cite{X.4}, inspir\'es par les
r\'esultats de \textsc{Zariski} \cite[chap.\, VIII]{X.6}, d\'emontr\'es par voie
topologique. Ces derniers sont loin d'avoir \'et\'e assimil\'es
par la g\'eom\'etrie alg\'ebrique \og abstraite\fg et m\'eritent
de nouveaux efforts.
\end{remarque}

Nous
aurons besoin de quelques faits \'el\'ementaires de la th\'eorie
de la ramification. Soient $V$ un anneau de valuation discr\`ete de
corps des fractions~$K$, corps r\'esiduel~$k$, $L$ une extension
galoisienne de~$K$ de groupe~$G$, $V'$ le normalis\'e de~$V$
dans~$L$, qui est un module libre de rang $n=[L{\colon }K]$ sur~$V$,
$\goth{m'}$ un id\'eal maximal de~$V'$, $G_d$ le sous-groupe de~$G$
form\'e des \'el\'ements laissant $\goth{m'}$ invariant, de
sorte que $G_d$ op\`ere dans l'extension r\'esiduelle
$k'=V'/\goth{m'}$ de~$k$, et $G_i$ le sous-groupe des \'el\'ements
de~$G_d$ op\'erant trivialement (rappelons que $G_d$ et $G_i$ sont
appel\'es respectivement sous-groupes de \emph{décomposition} et
d'\emph{inertie} de~$G$). On dit que $L$ est \emph{modérément
ramifié} sur~$V$ si $n_i=[G_i{\colon }e]$ est d'ordre $p$ premier
\`a la caract\'eristique de~$k$ (condition toujours
v\'erifi\'ee si $k$ est de caract\'eristique~0). Il est bien
connu que $G_i$ se plonge alors canoniquement dans le groupe ${k'}^*$,
donc est isomorphe au groupe des racines
$n_i$-i\`emes de l'unit\'e
dans~${k'}^*$, ce qui implique en particulier que $G_i$ \emph{est
cyclique}. Le cas type de cette situation est celui o\`u on pose
$L=K[t]/(t^n - u)$, $u$ \'etant une uniformisante de~$V$ et $n$ un
entier premier \`a~$p$: si $K$ contient les racines \niemes de
l'unit\'e, $L$ est une extension galoisienne totalement ramifi\'ee
de~$K$, de groupe de Galois $G=G_i$ isomorphe \`a~$\ZZ/n\ZZ$.

\begin{lemme}[Lemme d'\textsc{Abhyankar}]
\label{X.3.6}
\index{Abhyankar (lemme d')|hyperpage}\index{lemme d'Abhyankar|hyperpage}Soit $V$ un anneau de valuation discr\`ete de
corps des fractions~$K$, $L$ et $K'$ deux extensions galoisiennes
de~$K$ \emph{modérément ramifiées} sur~$V$, $n$ et $m$ les
ordres des groupes d'inertie
correspondants,
$L'$ une extension
compos\'ee de~$L$ et $K'$ sur~$K$. Si $m$ est un multiple de~$n$,
alors $L'$ est non ramifi\'ee sur les localis\'es de la
cl\^oture normale $V'$ de~$V$ dans~$K'$.
\end{lemme}

Soient en effet $W'$ le normalis\'e de~$V'$ dans~$L'$, $\goth{m'}$
un id\'eal maximal de~$V'$, $\goth{n'}$ un id\'eal maximal de~$W'$
au-dessus de
$\goth{m'}$,
$\goth{n}$ l'id\'eal maximal qu'il induit sur le normalis\'e $W$
de~$V$ dans~$L$, $G,\ H,\ M$ les groupes de Galois de~$L,\ K',\ L'$
sur~$K$, et $G_i,\ H_i,\ L_i'$ les groupes d'inertie correspondant aux
id\'eaux maximaux choisis. Alors $M$ se plonge dans le produit
$G\times H$ et $M_i$ dans le produit $G_i\times H_i$, de telle fa\c con que les projections $M\to G$ et $M\to H$, $M_i\to G_i$
et $M_i\to H_i$ soient surjectives (sorite du corps
interm\'ediaire). Il en r\'esulte d\'ej\`a, puisque $G_i$ et
$H_i$ sont par hypoth\`ese cycliques d'ordres $m$ et $n$ premiers
\`a~$p$, que $M_i$ est d'ordre premier \`a~$p$, donc cyclique, et
comme $m$ est multiple de~$n$ donc les \'el\'ements de~$G_i\times
H_i$ sont de puissance
$m$-i\`eme nulle, $M_i$ est d'ordre
divisant~$m$, donc d'ordre \'egal \`a $m$ puisque $M_i\to H_i$ est
surjectif. Ce dernier homomorphisme est donc \'egalement
injectif. Or son noyau est le groupe d'inertie de~$\goth{n'}$
au-dessus de~$\goth{m'}$, ce qui prouve que $L'$ est non ramifi\'e
sur~$K'$ en~$\goth{n'}$. D'o\`u le lemme.

Pla\c cons-nous maintenant sous les conditions de~\ref{X.2.4},
o\`u on a un homomorphisme de sp\'ecialisation
$$
\pi_1(\overline X_1,\overline a_1)\to\pi_1(\overline X_0,\overline
a_0)
$$
qui est \emph{surjectif}, relativement \`a un morphisme propre et
s\'eparable $f\colon X\to Y$. Nous voulons pr\'eciser le noyau de
cet homomorphisme. Proc\'edant comme dans la d\'emonstration
de~\ref{X.2.4}, on voit que dans cette question, on peut toujours
supposer que $Y$ est le spectre d'un \emph{anneau de valuation
discrète~$V$, complet et à corps résiduel
algébriquement clos} (car on peut toujours trouver un tel anneau
et un morphisme de son spectre $Y'$ dans $Y$ dont l'image soit
$\{y_0,y_1\}$). Alors on a $X_0=\overline X_0$, $\kres(y_0)=k=$ corps
r\'esiduel de~$V$, $\kres(y_1)=K=$ corps des fractions de~$V$. Soit
$K_s$ la cl\^oture s\'eparable de~$K$, $\overline K$ sa
cl\^oture alg\'ebrique, et pour tout sous-anneau $W$ de~$\overline
K$ contenant~$V$, posons $X_W=X\otimes_V W$. En particulier on a
$$
X_V=X, \qquad X_K=X_1, \qquad X_{\overline K}=\overline X_1.
$$
D'ailleurs le morphisme canonique $\overline X_1=X_{\overline K}\to
X_{K_s}$ induit un isomorphisme sur les groupes fondamentaux
(IX~\ref{IX.4.11}) de sorte que, compte tenu de
l'isomorphisme~\ref{X.2.1} $\pi_1(X_0)\to\pi_1(X)$, on est ramen\'e
\`a \'etudier l'homomorphisme surjectif
$$
\pi_1(X_{K_s})\to\pi_1(X)
$$
associ\'e
au morphisme canonique $X_{K_s}\to X$. La d\'etermination du noyau
de ce dernier revient \`a la solution du probl\`eme suivant:
\emph{on a un revêtement principal connexe $Z_{K_s}$ de
$X_{K_s}$, de groupe~$G$} (donc associ\'e \`a un homomorphisme de
$\pi_1(X_{K_s})$ dans~$G$), \emph{déterminer sous quelles
conditions il est isomorphe à l'image réciproque d'un
revêtement principal $Z$ de~$X$ de groupe~$G$.}

Notons d'abord que $K_s$ est r\'eunion filtrante croissante de ses
sous-extensions finies $K'$ sur~$K$, et que par suite, $Z_{K_s}$
est isomorphe \`a l'image inverse d'un rev\^etement principal
$Z_{K'}$ de~$X_{K'}$ pour un $K'$ convenable (on fera attention
cependant que pour $K'$ fix\'e, $Z_{K'}$ n'est pas d\'etermin\'e
de fa\c con unique). Dire que $Z_{K_s}$ est isomorphe \`a l'image
inverse d'un rev\^etement principal $Z$ de~$X$ signifie qu'il existe
une sous-extension finie $K''\supset K'$ de~$K_s$ telle que
$Z_{K''}=Z_{K'}\otimes_{K'} K''$ soit isomorphe \`a $Z\otimes_V
K''$. D\'esignons maintenant pour une sous-extension finie $K'$
de~$K_s$, par $V'$ le normalis\'e de~$V$ dans~$K'$, qui est un
anneau de valuation discr\`ete, complet, de corps
r\'esiduel~$k$. Donc le morphisme canonique $X_{V'}\to X_V$ induit
un isomorphisme pour les fibres au-dessus des points ferm\'es de
$Y=\Spec(V)$ et $Y'=\Spec(V')$ et il r\'esulte alors de~\ref{X.2.1}
appliqu\'e \`a $X_V$ et $X_{V'}$ que l'homomorphisme induit pour
les groupes fondamentaux $\pi_1(X_{V'})\to \pi_1(X_V)$ est un
isomorphisme, ou encore que tout rev\^etement principal de~$X_{V'}$
est l'image inverse d'un rev\^etement principal de~$X_V$
d\'etermin\'e \`a isomorphisme pr\`es. Cela implique donc le

\begin{lemme}
\label{X.3.7}
Soit $Z_{K'}$ un rev\^etement principal connexe de~$X_{K'}$ de
groupe~$G$, $Z_{K}$ son image inverse sur~$X_{K_s}$\quoi. Pour que ce
dernier soit isomorphe \`a l'image inverse d'un rev\^etement
principal $Z$ de~$X$, il faut et il suffit qu'il existe une extension
finie $K''\supset K'$ de~$K$ dans $K_s$ telle que le rev\^etement
principal $Z_{K''}$ de~$X_{K''}$ soit induit par un rev\^etement
principal de~$X_{V''}$.
\end{lemme}

Supposons en particulier que les $X_{V''}$ soient normaux (il suffit
par exemple pour cela que $X_0$ soit normal, et a fortiori que $X_0$
soit simple, \cf I~\ref{I.9.1}).
Comme ils sont connexes, ils sont donc irr\'eductibles. Soient $L$ le
corps des fonctions rationnelles pour $X$ et $X_{K}$, $L'$ celui pour
$X_{V'}$ et $X_{K'}$, $L''$ celui pour $X_{V''}$ et~$X_{K''}$\quoi. Alors
sous les conditions de~\ref{X.3.7}, $Z_{K'}$ d\'efinit une extension
finie s\'eparable $R'$ de~$L'$, et $Z_{K''}$ d\'efinit l'extension
$R''=R'\otimes_{L'} L'' = R'\otimes_{K'} K''$. La condition
envisag\'ee dans~\ref{X.3.7} signifie donc aussi qu'il existe une
extension finie s\'eparable $K''$ de~$K'$ telle que
$R''=R'\otimes_{K'} K''$ soit \emph{non ramifiée} au-dessus du
sch\'ema normal $X_{V''}$ de corps $L''=L'\otimes_{K'} K''$, et non
seulement au-dessus de la partie ouverte $X_{K''}$ de~$X_{V''}$.

Nous supposons dor\'enavant que $f\colon X\to Y$ est un morphisme
\emph{lisse,} donc les morphismes $X_{V'}\to \Spec(V')$ sont lisses,
donc les sch\'emas $X_{V'}$ sont \emph{réguliers.} Noter que la
fibre du point ferm\'e de~$\Spec(V')$ dans $X_{V'}$ est
irr\'eductible et de codimension~$1$. Soit $\goth{o'}$ son anneau
local, qui est donc un anneau de valuation discr\`ete de corps
$L'$, de corps r\'esiduel isomorphe au corps des fonctions
rationnelles de~$X_0$, donc ayant m\^eme caract\'eristique que
$k$. D\'efinissons de m\^eme $\goth{o''}$ dans $L''$, qui est
\'evidemment le normalis\'e de
$\goth{o'}$
dans~$L''$. Il
r\'esulte alors du th\'eor\`eme de puret\'e~\ref{X.3.1}
ou~\ref{X.3.3} que pour que $R''$ soit non ramifi\'e sur~$X_{V''}$,
il faut et il suffit que $R''$ soit non ramifi\'e sur~$\goth{o''}$,
normalis\'e de~$\goth{o'}$ dans~$L''$.

Notons maintenant que si $u'$ est une uniformisante de~$V'$, c'est
aussi une uniformisante de~$\goth{o'}$. Si alors $n$ est un entier
premier \`a la caract\'eristique $p$ de~$k$, et si on prend
$K''=K'[t]/(t^n-u)$, alors
$K''$
est une extension galoisienne finie
de~$K'$ et $L''$ est isomorphe \`a $L'[t]/(t^n-u')$, donc est
mod\'er\'ement ramifi\'e sur~$\goth{o'}$ et de groupe d'inertie
d'ordre~$n$. Supposons alors $G$ \emph{d'ordre premier à~$p$,} ce
qui implique que $R'$ est mod\'er\'ement ramifi\'e sur~$\goth{o'}$, et prenons pour $n$ un multiple premier \`a $p$ de
l'ordre du groupe d'inertie de~$R'$ sur~$\goth{o'}$ (par exemple
$n=[G\colon e]$). Appliquant le
lemme d'\textsc{Abhyankar}~\ref{X.3.6}, on
voit que la condition envisag\'ee dans~\ref{X.3.7} est
v\'erifi\'ee.

Cela prouve le th\'eor\`eme suivant:

\begin{theoreme}
\label{X.3.8}
Soient $f\colon X\to Y$ un morphisme propre et lisse, \`a fibres
g\'eo\-m\'e\-tri\-que\-ment connexes, avec $Y$ localement noeth\'erien,
$y_0$ et $y_1$ deux points de~$Y$ tels que
$y_0\in\overline{\{y_1\}}$,
$\overline{X}_0$ et $\overline{X}_1$ les fibres g\'eom\'etriques
correspondantes, consid\'erons l'homomorphisme de sp\'ecialisation
\eqref{X.2.4} $\pi_1(\overline X_1)\to\pi_1(\overline X_0)$. Cet
homomorphisme est surjectif, et tout homomorphisme continu de
$\pi_1(\overline X_1)$ dans un groupe fini $G$ d'ordre premier \`a
la caract\'eristique $p$ de~$\kres(y_0)$ provient d'un homomorphisme de
$\pi_1(\overline{X}_0)$ dans~$G$.
\end{theoreme}

En d'autres termes:

\begin{corollaire}
\label{X.3.9}
Si $\kres(y_0)$ est de caract\'eristique nulle, alors l'homomorphisme
de sp\'ecialisation est un isomorphisme. Si
$\kres(y_0)$ est de caract\'eristique $p>0$,
alors le noyau de
l'homomorphisme de sp\'ecialisation est contenu dans l'intersection
des noyaux des homomorphismes continus de~$\pi_1(\overline X_1)$ dans
des groupes finis d'ordre premier \`a~$p$ (ou encore, le sous-groupe
invariant ferm\'e engendr\'e par un $p$-sous-groupe de Sylow du
groupe de type galoisien $\pi_1(\overline X_1)$); si donc
$\pi_1(\overline X_1)^{(p)}$
d\'esigne le groupe quotient de
$\pi_1(\overline X_1)$ par le sous-groupe ferm\'e pr\'ec\'edent,
et si on d\'efinit de m\^eme $\pi_1(\overline X_0)^{(p)}$, alors
l'homomorphisme de sp\'ecialisation induit un \emph{isomorphisme}
$$
\pi_1(\overline X_1)^{(p)}\isomto \pi_1(\overline X_0)^{(p)}
$$
\end{corollaire}

On notera que la d\'emonstration de \ref{X.3.8} est purement alg\'ebrique.
Proc\'edant comme dans \ref{X.2.6}, on en conclut \emph{par voie
transcendante:}

\begin{corollaire}
\label{X.3.10}
Soit $X_0$ une courbe propre, lisse et connexe de genre $g$ sur un
corps al\-g\'e\-bri\-quement clos de caract\'eristique~$p$. Avec
la notation introduite dans~\ref{X.3.9}, le groupe $\pi_1(X_0)^{(p)}$
est isomorphe \`a $\Gamma^{(p)}$, o\`u $\Gamma$ est le groupe de
type galoisien engendr\'e par des g\'en\'erateurs $s_i,~t_i$
($1\le i\le g$) li\'es par la relation
$$
(s_1t_1s_1^{-1}t_1^{-1})\cdots (s_gt_gs_g^{-1}t_g^{-1})=1
$$
\end{corollaire}

\begin{remarques}
\label{X.3.11}
Dans le cas o\`u $\kres(y_0)$ est de caract\'eristique nulle, le
r\'esultat~\ref{X.3.9} est bien connu par voie transcendante. On
notera que la d\'emonstration de~\ref{X.3.10} fait appel au
th\'eor\`eme de puret\'e dans le cas d'in\'egales
caract\'eristiques, mais dans le cas d'anneaux de dimension 2
seulement, o\`u la d\'emonstration dudit th\'eor\`eme est
facile et a \'et\'e rappel\'ee dans le texte.
\end{remarques}

\begin{thebibliography}{0}{X.4}

\bibitem[1]{X.1} K.\, \textsc{Kodaira}, On compact analytic surfaces,
Ann.\ of Math.\ \textbf{71} (1960), p.\, 111--152.

\bibitem[2]{X.2} S.\, \textsc{Lang} et
J-P.\,\textsc{Serre}, Sur les rev\^etements non
ramifi\'es des vari\'et\'es alg\'ebriques,
Amer.\ J.\ Math.\ \textbf{79} (1957), p.\,319--330.

\bibitem[3]{X.3} M.\, \textsc{Nagata},
On the purity of branch loci in regular local rings,
Illinois J.\ Math. \textbf{3} (1959), p.\, 328--333.

\bibitem[4]{X.4}
J-P.\,\textsc{Serre}, Rev\^etements ramifi\'es du plan
projectif (d'apr\`es S.\, Abhyankar), S\'eminaire Bourbaki, Mai
1960.

\bibitem[5]{X.5} O.\, \textsc{Zariski},
On the purity of the branch locus of algebraic functions,
Proc.\ Nat.\ Acad.\ Sc.\ USA  \textbf{44} (1958), p.\, 791--796.


\bibitem[6]{X.6} O.\, \textsc{Zariski},
Algebraic surfaces, Ergebnisse... 1948, Chelsea (New York).
\end{thebibliography}

\chapterspace{-2}
